\documentclass[10pt,a4paper]{amsart}
\usepackage{amsmath,amssymb,amsthm,graphicx}
\usepackage{verbatim}
\newenvironment{reference}{\comment}{\endcomment}
%\newenvironment{reference}{}{}
\newenvironment{slogan}{\comment}{\endcomment}
\newenvironment{history}{\comment}{\endcomment}
% For commutative diagrams we use Xy-pic
\usepackage[all]{xy}

% We use 2cell for 2-commutative diagrams.
\xyoption{2cell}
% Theorem environments.
%
\theoremstyle{plain}
\newtheorem{theorem}[subsection]{Theorem}
\newtheorem{proposition}[subsection]{Proposition}
\newtheorem{lemma}[subsection]{Lemma}

\theoremstyle{definition}
\newtheorem{definition}[subsection]{Definition}
\newtheorem{example}[subsection]{Example}
\newtheorem{exercise}[subsection]{Exercise}
\newtheorem{situation}[subsection]{Situation}

\theoremstyle{remark}
\newtheorem{remark}[subsection]{Remark}
\newtheorem{remarks}[subsection]{Remarks}

\numberwithin{equation}{subsection}

% Macros
%
\def\lim{\mathop{\mathrm{lim}}\nolimits}
\def\colim{\mathop{\mathrm{colim}}\nolimits}
\def\Spec{\mathop{\mathrm{Spec}}}
\def\Hom{\mathop{\mathrm{Hom}}\nolimits}
\def\Ext{\mathop{\mathrm{Ext}}\nolimits}
\def\SheafHom{\mathop{\mathcal{H}\!\mathit{om}}\nolimits}
\def\SheafExt{\mathop{\mathcal{E}\!\mathit{xt}}\nolimits}
\def\Sch{\mathit{Sch}}
\def\Mor{\mathop{\mathrm{Mor}}\nolimits}
\def\Ob{\mathop{\mathrm{Ob}}\nolimits}
\def\Sh{\mathop{\mathit{Sh}}\nolimits}
\def\NL{\mathop{N\!L}\nolimits}
\def\CH{\mathop{\mathrm{CH}}\nolimits}
\def\proetale{{pro\text{-}\acute{e}tale}}
\def\etale{{\acute{e}tale}}
\def\QCoh{\mathit{QCoh}}
\def\Ker{\mathop{\mathrm{Ker}}}
\def\Im{\mathop{\mathrm{Im}}}
\def\Coker{\mathop{\mathrm{Coker}}}
\def\Coim{\mathop{\mathrm{Coim}}}

% Boxtimes
%
\DeclareMathSymbol{\boxtimes}{\mathbin}{AMSa}{"02}

%
% Macros for moduli stacks/spaces
%
\def\QCohstack{\mathcal{QC}\!\mathit{oh}}
\def\Cohstack{\mathcal{C}\!\mathit{oh}}
\def\Spacesstack{\mathcal{S}\!\mathit{paces}}
\def\Quotfunctor{\mathrm{Quot}}
\def\Hilbfunctor{\mathrm{Hilb}}
\def\Curvesstack{\mathcal{C}\!\mathit{urves}}
\def\Polarizedstack{\mathcal{P}\!\mathit{olarized}}
\def\Complexesstack{\mathcal{C}\!\mathit{omplexes}}
% \Pic is the operator that assigns to X its picard group, usage \Pic(X)
% \Picardstack_{X/B} denotes the Picard stack of X over B
% \Picardfunctor_{X/B} denotes the Picard functor of X over B
\def\Pic{\mathop{\mathrm{Pic}}\nolimits}
\def\Picardstack{\mathcal{P}\!\mathit{ic}}
\def\Picardfunctor{\mathrm{Pic}}
\def\Deformationcategory{\mathcal{D}\!\mathit{ef}}
\usepackage{cleveref}
\crefname{originallemma}{original lemma}{original lemmas}
\Crefname{originallemma}{Original lemma}{Original lemmas}
\makeatletter
\newlabel{lemma-sober-subspace}{{lemma-sober-subspace}{source}{}{sourcepointer-1}{}}
\newlabel{lemma-sober-subspace@cref}{{[originallemma][0][]lemma-sober-subspace}{[1][0][]source}{}{}{}}
\makeatother
\title{Sobriety is local on an open covering with noetherian intersections}
\author{The Stacks Project Authors}
\date{stacks-project:topology.tex:local-snapshot-20261012; GNU FDL 1.2 or later}
\begin{document}
\maketitle
\noindent\textit{Editorial note.} All delivered mathematical text is the author's own, copied from the pinned Stacks Project chapter source sources/stacks-topology/pinned-original-source.tex; no statement or proof step was rewritten, reordered or removed.\par\medskip
\noindent\textit{Editorial note.} Preamble provenance: the Stacks Project keeps one single shared preamble for every chapter; this card takes it directly from the staged sources/stacks-topology/preamble.tex as three verbatim spans (lines 8-12, 14-18 and 154-217), so no derived file was written and nothing was added or redefined. The wrapper class and the base packages are supplied by the assembler.\par\medskip
\noindent\textit{Editorial note.} Bibliography provenance: neither the counted statement nor its proof contains a citation (this lane rescans the whole span before writing the card), so the card carries no bibliography and no entry was reconstructed.\par\medskip
\noindent\textit{Editorial note.} Selection provenance: this card was opened under the file-based control contract (CONTROL.md), with the mechanical pre-card gate run on its statement and proof interval before the card was written.\par\medskip
\noindent\textit{Allowed standard prerequisites (editorial).}\begin{itemize}
\item The characterization of sobriety through generic points of irreducible closed subsets. Used at the start of the counted proof (source lines 1262-1270). Reduces the assertion to producing generic points.
\end{itemize}\medskip
\ifdefined\hypertarget\hypertarget{sourcepointer-1}{}\fi
\noindent\textit{Source reference.} Original lemma lemma-sober-subspace: sources/stacks-topology/pinned-original-source.tex, source lines 1213--1213. Navigation pointer only: the counted text refers to the Stacks item lemma-sober-subspace, which is an existing labeled result of this same chapter and is not part of this card. The pointer carries the target's own label as its number and its definition line as its span; its content is not reproduced here.\par\medskip
\begin{lemma}
\label{lemma-sober-local}
Let $X$ be a topological space and let $(X_i)_{i\in I}$ be a covering of $X$.
\begin{enumerate}
\item Suppose $X_i$ is locally closed in $X$ for every $i\in I$. Then, $X$ is
Kolmogorov if and only if $X_i$ is Kolmogorov for every $i\in I$.
\item Suppose $X_i$ is open in $X$ for every $i\in I$. Then, $X$ is
quasi-sober if and only if $X_i$ is quasi-sober for every $i\in I$.
\item Suppose $X_i$ is open in $X$ for every $i\in I$. Then, $X$ is sober if
and  only if $X_i$ is sober for every $i\in I$.
\end{enumerate}
\end{lemma}

\begin{proof}
Proof of (1). If $X$ is Kolmogorov then so is $X_i$ for every $i\in I$ by
Lemma \ref{lemma-sober-subspace}. Suppose $X_i$ is Kolmogorov for every
$i\in I$. Let $x,y\in X$ with $\overline{\{x\}}=\overline{\{y\}}$. There exists
$i\in I$ with $x\in X_i$. There exists an open subset $U\subset X$ such that
$X_i$ is a closed subset of $U$. If $y\notin U$ we get the contradiction
$x\in\overline{\{x\}}\cap U=\overline{\{y\}}\cap U=\emptyset$. Hence $y\in U$.
It follows $y\in\overline{\{y\}}\cap U=\overline{\{x\}}\cap U\subset X_i$.
This shows $y\in X_i$. It follows
$\overline{\{x\}}\cap X_i=\overline{\{y\}}\cap X_i$. Since $X_i$ is
Kolmogorov we get $x=y$. This shows $X$ is Kolmogorov.

\medskip\noindent
Proof of (2). If $X$ is quasi-sober then so is $X_i$ for every $i\in I$ by
Lemma \ref{lemma-sober-subspace}. Suppose $X_i$ is quasi-sober for every
$i\in I$. Let $Y$ be an irreducible closed subset of $X$. As $Y\neq\emptyset$
there exists $i\in I$ with $X_i\cap Y\neq\emptyset$. As $X_i$ is open in $X$ it
follows $X_i\cap Y$ is non-empty and open in $Y$, hence irreducible
and dense in $Y$. Thus $X_i\cap Y$ is an irreducible closed subset of $X_i$. As
$X_i$ is quasi-sober there exists $x\in X_i\cap Y$ with
$X_i\cap Y=\overline{\{x\}}\cap X_i\subset\overline{\{x\}}$. Since
$X_i\cap Y$ is dense in $Y$ and $Y$ is closed in $X$ it follows
$Y=\overline{X_i\cap Y}\cap Y\subset\overline{X_i\cap Y}\subset
\overline{\{x\}}\subset Y$. Therefore
$Y=\overline{\{x\}}$. This shows $X$ is quasi-sober.

\medskip\noindent
Proof of (3). Immediately from (1) and (2).
\end{proof}

\end{document}
